% ============================================================================
% UNAL SPACE — Documentación Técnica de la Base de Datos
% Módulo de Espacios Deportivos · Sede Bogotá
% Para compilar en Overleaf con pdfLaTeX o XeLaTeX
% ============================================================================
\documentclass[12pt, a4paper]{article}

% ── Paquetes ────────────────────────────────────────────────────────────────
\usepackage[utf8]{inputenc}
\usepackage[T1]{fontenc}
\usepackage[spanish]{babel}
\usepackage{geometry}
\geometry{margin=2.5cm}
\usepackage{graphicx}
\usepackage{booktabs}
\usepackage{longtable}
\usepackage{array}
\usepackage{tabularx}
\usepackage{multirow}
\usepackage{xcolor}
\usepackage{colortbl}
\usepackage{listings}
\usepackage{hyperref}
\usepackage{enumitem}
\usepackage{tikz}
\usetikzlibrary{shapes.multipart, arrows.meta, positioning, calc, fit}
\usepackage{fancyhdr}
\usepackage{titlesec}
\usepackage{tocloft}
\usepackage{caption}
\usepackage{float}
\usepackage{pdflscape}
\usepackage{adjustbox}
\usepackage{amsmath}

% ── Colores institucionales UNAL ────────────────────────────────────────────
\definecolor{unalverde}{HTML}{006633}
\definecolor{unalgris}{HTML}{4D4D4D}
\definecolor{unalclaro}{HTML}{E8F5E9}
\definecolor{sqlblue}{HTML}{1565C0}
\definecolor{sqlgreen}{HTML}{2E7D32}
\definecolor{sqlred}{HTML}{C62828}
\definecolor{sqlorange}{HTML}{E65100}
\definecolor{codebg}{HTML}{F5F5F5}
\definecolor{codeframe}{HTML}{BDBDBD}

% ── Configuración de listings para SQL ──────────────────────────────────────
\lstdefinestyle{sqlstyle}{
    language=SQL,
    backgroundcolor=\color{codebg},
    basicstyle=\ttfamily\footnotesize,
    keywordstyle=\color{sqlblue}\bfseries,
    stringstyle=\color{sqlred},
    commentstyle=\color{sqlgreen}\itshape,
    morekeywords={ENUM, UUID, TIMESTAMPTZ, TSTZRANGE, BOOLEAN, VARCHAR, TEXT,
                  SERIAL, GENERATED, ALWAYS, STORED, EXCLUDE, USING, gist,
                  ISODOW, EPOCH, STABLE, LATERAL, EXTENSION, CONFLICT,
                  NOTHING, RESTRICT, CASCADE, BEFORE, AFTER, EACH, ROW,
                  EXECUTE, FUNCTION, TRIGGER, RETURNS, LANGUAGE, plpgsql,
                  DECLARE, BEGIN, END, IF, ELSIF, THEN, RAISE, EXCEPTION,
                  INTO, NEW, OLD, TG_OP, RETURN, INTERVAL, ZONE, TIME,
                  WITH, OVER, PARTITION, RANK, NULLIF, COALESCE, DISTINCT,
                  tstzrange, generate_series, uuid_generate_v4, NOW, EXTRACT,
                  AT, BETWEEN},
    morecomment=[l]{--},
    morestring=[b]',
    numbers=left,
    numberstyle=\tiny\color{gray},
    numbersep=8pt,
    frame=single,
    rulecolor=\color{codeframe},
    breaklines=true,
    breakatwhitespace=true,
    showstringspaces=false,
    tabsize=4,
    captionpos=b,
    xleftmargin=15pt,
    framexleftmargin=15pt,
    aboveskip=12pt,
    belowskip=8pt,
}
\lstset{style=sqlstyle}

% ── Encabezados y pies de página ────────────────────────────────────────────
\pagestyle{fancy}
\fancyhf{}
\fancyhead[L]{\footnotesize\textcolor{unalgris}{UNAL Space --- Documentación de Base de Datos}}
\fancyhead[R]{\footnotesize\textcolor{unalgris}{Universidad Nacional de Colombia}}
\fancyfoot[C]{\thepage}
\renewcommand{\headrulewidth}{0.4pt}
\renewcommand{\footrulewidth}{0.2pt}

% ── Formato de secciones ────────────────────────────────────────────────────
\titleformat{\section}
    {\Large\bfseries\color{unalverde}}{\thesection.}{0.5em}{}
\titleformat{\subsection}
    {\large\bfseries\color{unalgris}}{\thesubsection.}{0.5em}{}
\titleformat{\subsubsection}
    {\normalsize\bfseries\color{unalgris}}{\thesubsubsection.}{0.5em}{}

% ── Hyperref ────────────────────────────────────────────────────────────────
\hypersetup{
    colorlinks=true,
    linkcolor=unalverde,
    urlcolor=sqlblue,
    citecolor=unalverde,
    pdftitle={UNAL Space - Documentación de Base de Datos},
    pdfauthor={Ingeniería de Software - UNAL},
}

% ── Macros ──────────────────────────────────────────────────────────────────
\newcommand{\pk}{\texttt{PK}}
\newcommand{\fk}{\texttt{FK}}
\newcommand{\uq}{\texttt{UQ}}
\newcommand{\nn}{\texttt{NN}}
\newcommand{\chk}{\texttt{CK}}
\newcommand{\tname}[1]{\texttt{\textbf{#1}}}
\newcommand{\cname}[1]{\texttt{#1}}

% ════════════════════════════════════════════════════════════════════════════
%                           INICIO DEL DOCUMENTO
% ════════════════════════════════════════════════════════════════════════════
\begin{document}

% ── Portada ─────────────────────────────────────────────────────────────────
\begin{titlepage}
    \centering
    \vspace*{2cm}

    {\Huge\bfseries\textcolor{unalverde}{UNAL Space}\par}
    \vspace{0.5cm}
    {\LARGE Plataforma de Reservas de Espacios Deportivos\par}
    \vspace{0.3cm}
    {\large Universidad Nacional de Colombia --- Sede Bogotá\par}

    \vspace{2cm}

    \rule{\textwidth}{1pt}
    \vspace{0.5cm}
    {\Large\bfseries Documentación Técnica de la Base de Datos\par}
    \vspace{0.2cm}
    {\large Módulo de Espacios Deportivos\par}
    \vspace{0.5cm}
    \rule{\textwidth}{1pt}

    \vspace{2cm}

    \begin{tabular}{ll}
        \textbf{Motor de BD:}      & PostgreSQL 16+ \\
        \textbf{Zona horaria:}     & America/Bogota (UTC-5) \\
        \textbf{Base de datos:}    & \texttt{unal\_space} \\
        \textbf{Herramienta IDE:}  & JetBrains DataGrip \\
        \textbf{Versión del documento:} & 1.0 \\
    \end{tabular}

    \vfill
    {\large Ingeniería de Software\par}
    {\large Facultad de Ingeniería\par}
    {\large\the\year\par}
\end{titlepage}

% ── Tabla de contenido ──────────────────────────────────────────────────────
\tableofcontents
\newpage


% ════════════════════════════════════════════════════════════════════════════
\section{Introducción}
% ════════════════════════════════════════════════════════════════════════════

El presente documento describe en detalle el diseño, la estructura y la implementación de la base de datos relacional del módulo de \textbf{Espacios Deportivos} de la plataforma \textbf{UNAL Space}, desarrollada para la Universidad Nacional de Colombia, Sede Bogotá.

El sistema permite a los estudiantes de la universidad reservar espacios deportivos del campus (canchas, estadio, escenarios) de forma controlada, garantizando:

\begin{itemize}[leftmargin=2em]
    \item Autenticación mediante correo institucional (\texttt{@unal.edu.co}).
    \item Cumplimiento de horarios de servicio (lunes a viernes, 06:00--21:00).
    \item Control de duración de reservas (mínimo 1 hora, máximo 2 horas).
    \item Prevención absoluta de solapamientos mediante restricciones ACID a nivel de motor.
    \item Trazabilidad completa a través de auditoría automatizada.
\end{itemize}

\subsection{Tecnologías utilizadas}

\begin{table}[H]
\centering
\begin{tabular}{lp{10cm}}
    \toprule
    \textbf{Componente} & \textbf{Tecnología} \\
    \midrule
    Motor de BD & PostgreSQL 16+ \\
    Extensiones & \texttt{uuid-ossp} (generación de UUID v4), \texttt{btree\_gist} (índices GiST para exclusión de rangos) \\
    Tipos de dato & UUID, TIMESTAMPTZ, TSTZRANGE, ENUM personalizados \\
    IDE de BD & JetBrains DataGrip \\
    Zona horaria & \texttt{America/Bogota} (UTC-5) \\
    \bottomrule
\end{tabular}
\caption{Stack tecnológico de la base de datos.}
\end{table}


% ════════════════════════════════════════════════════════════════════════════
\section{Tipos Enumerados (ENUM)}
% ════════════════════════════════════════════════════════════════════════════

Se definen cuatro tipos enumerados que restringen los valores posibles en columnas clave, garantizando integridad de dominio a nivel de motor:

\begin{table}[H]
\centering
\renewcommand{\arraystretch}{1.3}
\begin{tabularx}{\textwidth}{l l X}
    \toprule
    \textbf{Tipo ENUM} & \textbf{Valores} & \textbf{Descripción} \\
    \midrule
    \texttt{tipo\_espacio} &
        \texttt{ESTADIO}, \texttt{CONCHA\_ACUSTICA}, \texttt{CANCHA\_TECHADA}, \texttt{CANCHA\_DESTECHADA}, \texttt{CANCHA\_TENIS} &
        Clasificación física del escenario deportivo. \\
    \addlinespace
    \texttt{tipo\_superficie} &
        \texttt{GRAMA\_NATURAL}, \texttt{GRAMA\_SINTETICA}, \texttt{CEMENTO}, \texttt{CAUCHO\_SINTETICO}, \texttt{POLVO\_DE\_LADRILLO}, \texttt{BALDOSA\_DEPORTIVA} &
        Material de la superficie de juego. \\
    \addlinespace
    \texttt{deporte\_permitido} &
        \texttt{FUTBOL\_11}, \texttt{FUTBOL\_SALA}, \texttt{BALONCESTO}, \texttt{VOLEIBOL}, \texttt{TENIS} &
        Disciplinas deportivas disponibles en el campus. \\
    \addlinespace
    \texttt{estado\_reserva} &
        \texttt{PENDIENTE}, \texttt{CONFIRMADA}, \texttt{CANCELADA}, \texttt{FINALIZADA}, \texttt{NO\_SHOW} &
        Ciclo de vida de una reserva. \\
    \bottomrule
\end{tabularx}
\caption{Tipos enumerados definidos en la base de datos.}
\end{table}

\begin{lstlisting}[caption={Definición de tipos enumerados en SQL.}]
CREATE TYPE tipo_espacio AS ENUM (
    'ESTADIO', 'CONCHA_ACUSTICA', 'CANCHA_TECHADA',
    'CANCHA_DESTECHADA', 'CANCHA_TENIS'
);

CREATE TYPE tipo_superficie AS ENUM (
    'GRAMA_NATURAL', 'GRAMA_SINTETICA', 'CEMENTO',
    'CAUCHO_SINTETICO', 'POLVO_DE_LADRILLO', 'BALDOSA_DEPORTIVA'
);

CREATE TYPE deporte_permitido AS ENUM (
    'FUTBOL_11', 'FUTBOL_SALA', 'BALONCESTO', 'VOLEIBOL', 'TENIS'
);

CREATE TYPE estado_reserva AS ENUM (
    'PENDIENTE', 'CONFIRMADA', 'CANCELADA', 'FINALIZADA', 'NO_SHOW'
);
\end{lstlisting}


% ════════════════════════════════════════════════════════════════════════════
\section{Modelo Entidad--Relación}
% ════════════════════════════════════════════════════════════════════════════

\subsection{Diagrama ER}

La Figura~\ref{fig:er-diagram} presenta el modelo entidad--relación completo del módulo deportivo. Las relaciones se describen con notación de cardinalidad UML.

\begin{figure}[H]
\centering
\resizebox{\textwidth}{!}{%
\begin{tikzpicture}[
    entity/.style={
        rectangle split,
        rectangle split parts=2,
        draw=unalverde,
        thick,
        fill=unalclaro,
        text width=5.5cm,
        font=\small,
        align=left,
        inner sep=6pt,
        rounded corners=3pt
    },
    relationship/.style={
        diamond,
        draw=sqlorange,
        thick,
        fill=sqlorange!10,
        aspect=2.5,
        font=\small\bfseries,
        inner sep=2pt,
    },
    arrow/.style={
        -Stealth,
        thick,
        color=unalgris,
    },
    cardinality/.style={
        font=\footnotesize\bfseries,
        color=sqlblue,
    },
]

% ── Entidades ──
\node[entity] (usuarios) {
    \textbf{usuarios}
    \nodepart{second}
    \underline{id} : UUID (PK) \\
    nombres : VARCHAR(100) \\
    apellidos : VARCHAR(100) \\
    documento : VARCHAR(20) [UQ] \\
    codigo\_estudiantil : VARCHAR(20) [UQ] \\
    correo : VARCHAR(150) [UQ] \\
    facultad : VARCHAR(120) \\
    activo : BOOLEAN
};

\node[entity, right=4.5cm of usuarios] (reservas) {
    \textbf{reservas}
    \nodepart{second}
    \underline{id} : UUID (PK) \\
    usuario\_id : UUID (FK) \\
    espacio\_id : UUID (FK) \\
    disciplina\_id : UUID (FK) \\
    fecha\_inicio : TIMESTAMPTZ \\
    fecha\_fin : TIMESTAMPTZ \\
    rango\_horario : TSTZRANGE \\
    estado : estado\_reserva
};

\node[entity, right=4.5cm of reservas] (espacios) {
    \textbf{espacios\_deportivos}
    \nodepart{second}
    \underline{id} : UUID (PK) \\
    nombre : VARCHAR(150) \\
    codigo\_interno : VARCHAR(30) [UQ] \\
    tipo : tipo\_espacio \\
    superficie : tipo\_superficie \\
    es\_techado : BOOLEAN \\
    capacidad : INT \\
    es\_mixta : BOOLEAN
};

\node[entity, below=3.5cm of reservas] (disciplinas) {
    \textbf{disciplinas\_deportivas}
    \nodepart{second}
    \underline{id} : UUID (PK) \\
    nombre : VARCHAR(60) \\
    codigo : deporte\_permitido [UQ] \\
    descripcion : TEXT
};

\node[entity, below=3.5cm of espacios] (espdeportes) {
    \textbf{espacio\_deportes}
    \nodepart{second}
    \underline{id} : UUID (PK) \\
    espacio\_id : UUID (FK) \\
    disciplina\_id : UUID (FK)
};

\node[entity, below=3.5cm of usuarios] (auditoria) {
    \textbf{auditoria\_reservas}
    \nodepart{second}
    \underline{id} : UUID (PK) \\
    reserva\_id : UUID (FK) \\
    estado\_anterior : estado\_reserva \\
    estado\_nuevo : estado\_reserva \\
    descripcion : TEXT \\
    realizado\_por : VARCHAR(150)
};

% ── Relaciones (líneas) ──
% usuarios 1 ──── 0..* reservas
\draw[arrow] (usuarios.east |- usuarios.second south) -- node[above, cardinality] {1} ++(0.8,0) -- node[above, cardinality] {0..*} (reservas.west |- reservas.second south);

% reservas *..1 ──── espacios
\draw[arrow] (reservas.east |- reservas.second south) -- node[above, cardinality] {0..*} ++(0.8,0) -- node[above, cardinality] {1} (espacios.west |- espacios.second south);

% reservas *..1 ──── disciplinas
\draw[arrow] (reservas.south) -- node[right, cardinality] {0..*} ++(0,-0.8) -- node[right, cardinality] {1} (disciplinas.north);

% reservas 1 ──── 0..* auditoria
\draw[arrow] (reservas.west |- reservas.south) -- node[left, cardinality, xshift=-3pt] {1} ++(-2.2,-1.5) -- node[left, cardinality, xshift=-3pt] {0..*} (auditoria.north east);

% espacios 1 ──── 0..* espacio_deportes
\draw[arrow] (espacios.south) -- node[right, cardinality] {1} ++(0,-0.8) -- node[right, cardinality] {0..*} (espdeportes.north);

% disciplinas 1 ──── 0..* espacio_deportes
\draw[arrow] (disciplinas.east) -- node[above, cardinality] {1} ++(1.5,0) -- node[above, cardinality] {0..*} (espdeportes.west);

\end{tikzpicture}
}%
\caption{Diagrama Entidad--Relación del módulo de espacios deportivos de UNAL Space.}
\label{fig:er-diagram}
\end{figure}


\subsection{Cardinalidades}

\begin{table}[H]
\centering
\renewcommand{\arraystretch}{1.4}
\begin{tabularx}{\textwidth}{l c l X}
    \toprule
    \textbf{Entidad A} & \textbf{Cardinalidad} & \textbf{Entidad B} & \textbf{Descripción} \\
    \midrule
    \tname{usuarios} & $1 \longleftrightarrow 0..*$ & \tname{reservas} &
        Un estudiante puede tener cero o muchas reservas. Cada reserva pertenece a exactamente un estudiante. \\
    \addlinespace
    \tname{espacios\_deportivos} & $1 \longleftrightarrow 0..*$ & \tname{reservas} &
        Un espacio deportivo puede tener cero o muchas reservas. Cada reserva se realiza sobre exactamente un espacio. \\
    \addlinespace
    \tname{disciplinas\_deportivas} & $1 \longleftrightarrow 0..*$ & \tname{reservas} &
        Una disciplina puede estar asociada a muchas reservas. Cada reserva especifica exactamente una disciplina. \\
    \addlinespace
    \tname{espacios\_deportivos} & $1 \longleftrightarrow 0..*$ & \tname{espacio\_deportes} &
        Un espacio admite uno o más deportes. \\
    \addlinespace
    \tname{disciplinas\_deportivas} & $1 \longleftrightarrow 0..*$ & \tname{espacio\_deportes} &
        Una disciplina puede practicarse en uno o más espacios. \\
    \addlinespace
    \tname{reservas} & $1 \longleftrightarrow 0..*$ & \tname{auditoria\_reservas} &
        Cada reserva puede tener múltiples registros de auditoría (uno por cada cambio de estado). \\
    \bottomrule
\end{tabularx}
\caption{Tabla de cardinalidades entre entidades.}
\end{table}

\textbf{Nota sobre la relación muchos--a--muchos:} La relación entre \tname{espacios\_deportivos} y \tname{disciplinas\_deportivas} es de tipo $M:N$ y se resuelve a través de la tabla intermedia \tname{espacio\_deportes}, que define qué deportes se pueden practicar en cada espacio.


% ════════════════════════════════════════════════════════════════════════════
\section{Diccionario de Datos}
% ════════════════════════════════════════════════════════════════════════════

A continuación se documenta cada tabla con sus columnas, tipos de dato, restricciones y descripción funcional.

% ────────────────────────────────────────────────────────────────────────────
\subsection{Tabla: \tname{usuarios}}

Almacena la información de los estudiantes autorizados para realizar reservas.

\begin{longtable}{l l c l p{5cm}}
    \toprule
    \textbf{Columna} & \textbf{Tipo} & \textbf{Restr.} & \textbf{Default} & \textbf{Descripción} \\
    \midrule
    \endhead
    \cname{id} & UUID & \pk & \texttt{uuid\_generate\_v4()} & Identificador único. \\
    \cname{nombres} & VARCHAR(100) & \nn & --- & Nombres del estudiante. \\
    \cname{apellidos} & VARCHAR(100) & \nn & --- & Apellidos del estudiante. \\
    \cname{documento} & VARCHAR(20) & \nn, \uq & --- & Cédula de ciudadanía o tarjeta de identidad. \\
    \cname{codigo\_estudiantil} & VARCHAR(20) & \nn, \uq & --- & Código de matrícula universitaria. \\
    \cname{correo} & VARCHAR(150) & \nn, \uq, \chk & --- & Correo institucional. Validado con regex: \texttt{@unal.edu.co}. \\
    \cname{facultad} & VARCHAR(120) & \nn & --- & Facultad a la que pertenece. \\
    \cname{activo} & BOOLEAN & \nn & \texttt{TRUE} & Indica si el usuario está habilitado. \\
    \cname{creado\_en} & TIMESTAMPTZ & \nn & \texttt{NOW()} & Fecha y hora de creación. \\
    \cname{actualizado\_en} & TIMESTAMPTZ & \nn & \texttt{NOW()} & Última modificación (auto-actualizada por trigger). \\
    \bottomrule
    \caption{Diccionario de datos --- \tname{usuarios}.}
\end{longtable}

\textbf{Constraint de correo institucional:}
\begin{lstlisting}[caption={Validación de correo @unal.edu.co.}]
CONSTRAINT ck_usuario_correo_unal
    CHECK (correo ~* '^[A-Za-z0-9._%+-]+@unal\.edu\.co$')
\end{lstlisting}

% ────────────────────────────────────────────────────────────────────────────
\subsection{Tabla: \tname{disciplinas\_deportivas}}

Catálogo maestro de disciplinas deportivas disponibles en el campus.

\begin{longtable}{l l c l p{5cm}}
    \toprule
    \textbf{Columna} & \textbf{Tipo} & \textbf{Restr.} & \textbf{Default} & \textbf{Descripción} \\
    \midrule
    \endhead
    \cname{id} & UUID & \pk & \texttt{uuid\_generate\_v4()} & Identificador único. \\
    \cname{nombre} & VARCHAR(60) & \nn & --- & Nombre legible de la disciplina. \\
    \cname{codigo} & \texttt{deporte\_permitido} & \nn, \uq & --- & Código ENUM de la disciplina. \\
    \cname{descripcion} & TEXT & --- & --- & Descripción opcional. \\
    \cname{activo} & BOOLEAN & \nn & \texttt{TRUE} & Estado operativo. \\
    \cname{creado\_en} & TIMESTAMPTZ & \nn & \texttt{NOW()} & Fecha de creación. \\
    \bottomrule
    \caption{Diccionario de datos --- \tname{disciplinas\_deportivas}.}
\end{longtable}

% ────────────────────────────────────────────────────────────────────────────
\subsection{Tabla: \tname{espacios\_deportivos}}

Inventario físico de canchas, estadios y escenarios deportivos del campus.

\begin{longtable}{l l c l p{4.5cm}}
    \toprule
    \textbf{Columna} & \textbf{Tipo} & \textbf{Restr.} & \textbf{Default} & \textbf{Descripción} \\
    \midrule
    \endhead
    \cname{id} & UUID & \pk & \texttt{uuid\_generate\_v4()} & Identificador único. \\
    \cname{nombre} & VARCHAR(150) & \nn & --- & Nombre oficial del espacio. \\
    \cname{codigo\_interno} & VARCHAR(30) & \nn, \uq & --- & Código corto (ej. \texttt{TCH-001}). \\
    \cname{descripcion} & TEXT & --- & --- & Descripción detallada. \\
    \cname{tipo} & \texttt{tipo\_espacio} & \nn & --- & Clasificación física. \\
    \cname{superficie} & \texttt{tipo\_superficie} & \nn & --- & Material de superficie. \\
    \cname{ubicacion} & VARCHAR(200) & --- & --- & Ubicación referencial en el campus. \\
    \cname{edificio} & VARCHAR(80) & --- & --- & Código/nombre de edificio si aplica. \\
    \cname{es\_techado} & BOOLEAN & \nn & \texttt{FALSE} & Si tiene techo o cubierta. \\
    \cname{capacidad} & INT & \nn, \chk & \texttt{0} & Aforo del espacio ($\geq 0$). \\
    \cname{es\_mixta} & BOOLEAN & \nn & \texttt{FALSE} & Si admite múltiples disciplinas. \\
    \cname{activo} & BOOLEAN & \nn & \texttt{TRUE} & Estado operativo. \\
    \cname{creado\_en} & TIMESTAMPTZ & \nn & \texttt{NOW()} & Fecha de creación. \\
    \cname{actualizado\_en} & TIMESTAMPTZ & \nn & \texttt{NOW()} & Última modificación. \\
    \bottomrule
    \caption{Diccionario de datos --- \tname{espacios\_deportivos}.}
\end{longtable}

% ────────────────────────────────────────────────────────────────────────────
\subsection{Tabla: \tname{espacio\_deportes} (tabla intermedia M:N)}

Resuelve la relación muchos--a--muchos entre espacios y disciplinas deportivas, definiendo qué deportes se pueden practicar en cada espacio.

\begin{longtable}{l l c l p{5cm}}
    \toprule
    \textbf{Columna} & \textbf{Tipo} & \textbf{Restr.} & \textbf{Default} & \textbf{Descripción} \\
    \midrule
    \endhead
    \cname{id} & UUID & \pk & \texttt{uuid\_generate\_v4()} & Identificador único. \\
    \cname{espacio\_id} & UUID & \fk, \nn & --- & Referencia a \tname{espacios\_deportivos}. \\
    \cname{disciplina\_id} & UUID & \fk, \nn & --- & Referencia a \tname{disciplinas\_deportivas}. \\
    \cname{creado\_en} & TIMESTAMPTZ & \nn & \texttt{NOW()} & Fecha de creación. \\
    \bottomrule
    \caption{Diccionario de datos --- \tname{espacio\_deportes}.}
\end{longtable}

\textbf{Constraint compuesto:} \texttt{UNIQUE (espacio\_id, disciplina\_id)} --- una disciplina solo se asocia una vez por espacio.

% ────────────────────────────────────────────────────────────────────────────
\subsection{Tabla: \tname{reservas}}

Tabla central del sistema. Almacena cada reserva de un espacio deportivo por un estudiante.

\begin{longtable}{l l c l p{4.2cm}}
    \toprule
    \textbf{Columna} & \textbf{Tipo} & \textbf{Restr.} & \textbf{Default} & \textbf{Descripción} \\
    \midrule
    \endhead
    \cname{id} & UUID & \pk & \texttt{uuid\_generate\_v4()} & Identificador único. \\
    \cname{usuario\_id} & UUID & \fk, \nn & --- & Estudiante que reserva. \\
    \cname{espacio\_id} & UUID & \fk, \nn & --- & Espacio reservado. \\
    \cname{disciplina\_id} & UUID & \fk, \nn & --- & Deporte a practicar. \\
    \cname{fecha\_inicio} & TIMESTAMPTZ & \nn & --- & Inicio de la reserva. \\
    \cname{fecha\_fin} & TIMESTAMPTZ & \nn & --- & Fin de la reserva. \\
    \cname{rango\_horario} & TSTZRANGE & \nn & \textit{GENERATED} & Rango temporal auto-generado (\texttt{[inicio, fin)}). \\
    \cname{estado} & \texttt{estado\_reserva} & \nn & \texttt{PENDIENTE} & Estado actual del ciclo de vida. \\
    \cname{motivo\_cancelacion} & TEXT & --- & --- & Razón de cancelación si aplica. \\
    \cname{creado\_en} & TIMESTAMPTZ & \nn & \texttt{NOW()} & Fecha de creación. \\
    \cname{actualizado\_en} & TIMESTAMPTZ & \nn & \texttt{NOW()} & Última modificación. \\
    \bottomrule
    \caption{Diccionario de datos --- \tname{reservas}.}
\end{longtable}

\textbf{Columna generada:}
\begin{lstlisting}[caption={Rango temporal generado automáticamente.}]
rango_horario TSTZRANGE NOT NULL
    GENERATED ALWAYS AS (tstzrange(fecha_inicio, fecha_fin, '[)')) STORED
\end{lstlisting}

% ────────────────────────────────────────────────────────────────────────────
\subsection{Tabla: \tname{auditoria\_reservas}}

Registro inmutable de cada cambio de estado en las reservas para trazabilidad completa.

\begin{longtable}{l l c l p{4.5cm}}
    \toprule
    \textbf{Columna} & \textbf{Tipo} & \textbf{Restr.} & \textbf{Default} & \textbf{Descripción} \\
    \midrule
    \endhead
    \cname{id} & UUID & \pk & \texttt{uuid\_generate\_v4()} & Identificador único. \\
    \cname{reserva\_id} & UUID & \fk, \nn & --- & Reserva auditada. \\
    \cname{estado\_anterior} & \texttt{estado\_reserva} & --- & --- & Estado previo (NULL en creación). \\
    \cname{estado\_nuevo} & \texttt{estado\_reserva} & \nn & --- & Nuevo estado. \\
    \cname{descripcion} & TEXT & --- & --- & Descripción del cambio. \\
    \cname{realizado\_por} & VARCHAR(150) & --- & --- & Quién ejecutó el cambio. \\
    \cname{creado\_en} & TIMESTAMPTZ & \nn & \texttt{NOW()} & Timestamp de auditoría. \\
    \bottomrule
    \caption{Diccionario de datos --- \tname{auditoria\_reservas}.}
\end{longtable}


% ════════════════════════════════════════════════════════════════════════════
\section{Constraints y Reglas de Negocio}
% ════════════════════════════════════════════════════════════════════════════

El sistema implementa una estrategia de \textbf{defensa en profundidad} con tres capas de validación:

\subsection{Capa 1: CHECK Constraints (nivel DDL)}

Restricciones declarativas evaluadas por el motor PostgreSQL de forma atómica:

\begin{table}[H]
\centering
\renewcommand{\arraystretch}{1.3}
\begin{tabularx}{\textwidth}{l l X}
    \toprule
    \textbf{Constraint} & \textbf{Tabla} & \textbf{Regla} \\
    \midrule
    \texttt{ck\_usuario\_correo\_unal} & \tname{usuarios} &
        El correo debe coincidir con el patrón \texttt{*@unal.edu.co} (regex case-insensitive). \\
    \texttt{ck\_reserva\_orden\_fechas} & \tname{reservas} &
        \texttt{fecha\_fin > fecha\_inicio}. \\
    \texttt{ck\_reserva\_duracion} & \tname{reservas} &
        Duración entre 3600 y 7200 segundos (1--2 horas). \\
    \texttt{ck\_reserva\_horario\_apertura} & \tname{reservas} &
        Hora local entre 06:00 y 21:00 (America/Bogota). \\
    \texttt{ck\_reserva\_dia\_habil} & \tname{reservas} &
        ISODOW entre 1 (lunes) y 5 (viernes). \\
    \texttt{ck\_espacio\_capacidad} & \tname{espacios\_deportivos} &
        \texttt{capacidad >= 0}. \\
    \bottomrule
\end{tabularx}
\caption{CHECK constraints implementados.}
\end{table}

\subsection{Capa 2: EXCLUDE USING GiST (anti-solapamiento ACID)}

Esta es la restricción más crítica del sistema. Utiliza la extensión \texttt{btree\_gist} y la columna generada \texttt{rango\_horario} (tipo \texttt{TSTZRANGE}) para garantizar a nivel de motor que \textbf{es físicamente imposible} crear dos reservas que se solapen en el mismo espacio:

\begin{lstlisting}[caption={Restricción de exclusión GiST para anti-solapamiento.}]
ALTER TABLE reservas
    ADD CONSTRAINT excl_reserva_sin_solapamiento
        EXCLUDE USING gist (
            espacio_id    WITH =,
            rango_horario WITH &&
        )
        WHERE (estado IN ('PENDIENTE', 'CONFIRMADA'));
\end{lstlisting}

\textbf{Semántica:} Si dos filas tienen el mismo \texttt{espacio\_id} (operador \texttt{=}) y sus rangos temporales se solapan (operador \texttt{\&\&}), y ambas están en estado activo (\texttt{PENDIENTE} o \texttt{CONFIRMADA}), PostgreSQL rechaza la inserción con un error de exclusión. Esta garantía es transaccional (ACID) y resistente a concurrencia.

\subsection{Capa 3: Triggers de Validación Programática}

Un trigger \texttt{BEFORE INSERT OR UPDATE} en la tabla \tname{reservas} ejecuta la función \texttt{fn\_validar\_reserva()} que realiza validaciones adicionales con mensajes de error descriptivos en español:

\begin{enumerate}[leftmargin=2em]
    \item \textbf{Día hábil:} Verifica que la fecha no sea sábado ni domingo.
    \item \textbf{Franja horaria:} Confirma que el inicio $\geq$ 06:00 y el fin $\leq$ 21:00.
    \item \textbf{Duración:} Valida que esté entre 60 y 120 minutos.
    \item \textbf{Solapamiento personal:} Verifica que el mismo estudiante no tenga otra reserva activa que se solape en la misma franja del mismo día.
    \item \textbf{Disciplina válida:} Confirma que la disciplina seleccionada esté habilitada para ese espacio (mediante la tabla \tname{espacio\_deportes}).
\end{enumerate}

\begin{lstlisting}[caption={Fragmento del trigger de validación de reservas.}]
CREATE OR REPLACE FUNCTION fn_validar_reserva()
RETURNS TRIGGER LANGUAGE plpgsql AS $$
DECLARE
    v_hora_inicio   TIME;
    v_hora_fin      TIME;
    v_dia_semana    INT;
    v_duracion_min  DOUBLE PRECISION;
    v_fecha_local   DATE;
    v_reservas_dia  INT;
BEGIN
    v_hora_inicio := (NEW.fecha_inicio AT TIME ZONE 'America/Bogota')::time;
    v_hora_fin    := (NEW.fecha_fin AT TIME ZONE 'America/Bogota')::time;
    v_dia_semana  := EXTRACT(ISODOW FROM
                        NEW.fecha_inicio AT TIME ZONE 'America/Bogota');
    v_duracion_min := EXTRACT(EPOCH FROM
                        (NEW.fecha_fin - NEW.fecha_inicio)) / 60.0;

    -- a) Solo dias habiles (lunes a viernes)
    IF v_dia_semana NOT BETWEEN 1 AND 5 THEN
        RAISE EXCEPTION '[UNAL SPACE] Solo Lunes a Viernes.';
    END IF;

    -- b) Franja horaria 06:00 - 21:00
    IF v_hora_inicio < '06:00:00'::time THEN
        RAISE EXCEPTION '[UNAL SPACE] Inicio anterior a 06:00.';
    END IF;
    IF v_hora_fin > '21:00:00'::time THEN
        RAISE EXCEPTION '[UNAL SPACE] Fin posterior a 21:00.';
    END IF;

    -- c) Duracion entre 60 y 120 minutos
    IF v_duracion_min < 60 OR v_duracion_min > 120 THEN
        RAISE EXCEPTION '[UNAL SPACE] Duracion fuera de rango.';
    END IF;

    -- d) Sin solapamiento personal del mismo estudiante
    SELECT COUNT(*) INTO v_reservas_dia FROM reservas r
    WHERE r.usuario_id = NEW.usuario_id
      AND r.estado IN ('PENDIENTE', 'CONFIRMADA')
      AND r.rango_horario && tstzrange(
            NEW.fecha_inicio, NEW.fecha_fin, '[)')
      AND (TG_OP = 'INSERT' OR r.id <> NEW.id);

    IF v_reservas_dia > 0 THEN
        RAISE EXCEPTION '[UNAL SPACE] Solapamiento personal.';
    END IF;

    NEW.actualizado_en := NOW();
    RETURN NEW;
END;
$$;
\end{lstlisting}


% ════════════════════════════════════════════════════════════════════════════
\section{Inventario de Espacios Deportivos}
% ════════════════════════════════════════════════════════════════════════════

El campus cuenta con \textbf{12 espacios deportivos} distribuidos en la zona deportiva:

\begin{landscape}
\begin{table}[H]
\centering
\renewcommand{\arraystretch}{1.2}
\small
\begin{tabular}{c l l l l c c l}
    \toprule
    \textbf{\#} & \textbf{Código} & \textbf{Nombre} & \textbf{Tipo} & \textbf{Superficie} & \textbf{Techado} & \textbf{Cap.} & \textbf{Deportes} \\
    \midrule
    1  & \texttt{EST-001} & Estadio Alfonso López Pumarejo        & Estadio          & Grama natural      & No & 6000 & Fútbol 11 \\
    2  & \texttt{CON-001} & Concha Acústica -- Cancha Mixta 1      & Concha Acústica  & Cemento            & Sí & 40   & F.~Sala, Basquet, Volei \\
    3  & \texttt{CON-002} & Concha Acústica -- Cancha Mixta 2      & Concha Acústica  & Cemento            & Sí & 40   & F.~Sala, Basquet, Volei \\
    4  & \texttt{TCH-001} & Cancha Techada Mixta 1                 & Techada          & Caucho sintético   & Sí & 30   & F.~Sala, Basquet, Volei \\
    5  & \texttt{TCH-002} & Cancha Techada Mixta 2                 & Techada          & Caucho sintético   & Sí & 30   & F.~Sala, Basquet, Volei \\
    6  & \texttt{DES-001} & Cancha Destechada Mixta 1              & Destechada       & Cemento            & No & 30   & F.~Sala, Basquet, Volei \\
    7  & \texttt{DES-002} & Cancha Destechada Mixta 2              & Destechada       & Cemento            & No & 30   & F.~Sala, Basquet, Volei \\
    8  & \texttt{DES-003} & Cancha Destechada Mixta 3              & Destechada       & Cemento            & No & 30   & F.~Sala, Basquet, Volei \\
    9  & \texttt{BAL-001} & Cancha de Baloncesto Azul              & Destechada       & Caucho sintético   & No & 20   & Baloncesto \\
    10 & \texttt{TEN-001} & Cancha de Tenis 1                      & Tenis            & Polvo de ladrillo  & No & 4    & Tenis \\
    11 & \texttt{TEN-002} & Cancha de Tenis 2                      & Tenis            & Polvo de ladrillo  & No & 4    & Tenis \\
    12 & \texttt{TEN-003} & Cancha de Tenis 3                      & Tenis            & Polvo de ladrillo  & No & 4    & Tenis \\
    \bottomrule
\end{tabular}
\caption{Inventario completo de los 12 espacios deportivos del campus.}
\label{tab:inventario}
\end{table}
\end{landscape}


% ════════════════════════════════════════════════════════════════════════════
\section{Sistema de Auditoría}
% ════════════════════════════════════════════════════════════════════════════

El sistema implementa un mecanismo de auditoría automatizada mediante un trigger \texttt{AFTER INSERT OR UPDATE} en la tabla \tname{reservas} que registra cada cambio de estado en la tabla \tname{auditoria\_reservas}:

\begin{lstlisting}[caption={Trigger de auditoría (AFTER --- ejecuta después de confirmar la fila).}]
CREATE OR REPLACE FUNCTION fn_auditar_estado_reserva()
RETURNS TRIGGER LANGUAGE plpgsql AS $$
BEGIN
    IF TG_OP = 'INSERT' THEN
        INSERT INTO auditoria_reservas
            (reserva_id, estado_anterior, estado_nuevo,
             descripcion, realizado_por)
        VALUES (NEW.id, NULL, NEW.estado,
                'Reserva creada.', 'SISTEMA');
    ELSIF TG_OP = 'UPDATE'
          AND OLD.estado IS DISTINCT FROM NEW.estado THEN
        INSERT INTO auditoria_reservas
            (reserva_id, estado_anterior, estado_nuevo,
             descripcion, realizado_por)
        VALUES (NEW.id, OLD.estado, NEW.estado,
            CASE NEW.estado
                WHEN 'CONFIRMADA' THEN 'Reserva confirmada.'
                WHEN 'CANCELADA'  THEN 'Reserva cancelada.'
                WHEN 'FINALIZADA' THEN 'Bloque horario transcurrido.'
                WHEN 'NO_SHOW'    THEN 'Estudiante no se presento.'
                ELSE 'Cambio de estado.'
            END, 'SISTEMA');
    END IF;
    RETURN NULL;  -- AFTER trigger: retorno ignorado
END;
$$;
\end{lstlisting}

\textbf{Nota técnica:} El trigger es de tipo \texttt{AFTER} (no \texttt{BEFORE}) porque necesita referenciar la fila ya insertada en \tname{reservas} mediante la foreign key de \tname{auditoria\_reservas}. Un trigger \texttt{BEFORE} causaría el error \texttt{23503} (violación de llave foránea) al intentar insertar en la tabla de auditoría antes de que la reserva exista.

\subsection{Diagrama de estados de una reserva}

\begin{figure}[H]
\centering
\begin{tikzpicture}[
    state/.style={
        rectangle, rounded corners=8pt, draw=unalverde, thick,
        fill=unalclaro, minimum width=2.8cm, minimum height=0.9cm,
        font=\small\bfseries, text=unalgris
    },
    trans/.style={-Stealth, thick, color=unalgris},
    label/.style={font=\scriptsize, color=sqlorange, fill=white, inner sep=2pt}
]
    \node[state] (pendiente)   at (0,0)     {PENDIENTE};
    \node[state] (confirmada)  at (5,0)     {CONFIRMADA};
    \node[state] (cancelada)   at (0,-2.5)  {CANCELADA};
    \node[state] (finalizada)  at (5,-2.5)  {FINALIZADA};
    \node[state] (noshow)      at (10,-2.5) {NO\_SHOW};

    \draw[trans] (pendiente) -- node[label, above] {confirmar} (confirmada);
    \draw[trans] (pendiente) -- node[label, left]  {cancelar}  (cancelada);
    \draw[trans] (confirmada) -- node[label, right] {finalizar} (finalizada);
    \draw[trans] (confirmada) -- node[label, above] {no show}   (noshow);
    \draw[trans] (confirmada.south west) -- node[label, below, sloped] {cancelar} (cancelada.north east);
\end{tikzpicture}
\caption{Diagrama de transición de estados de una reserva.}
\label{fig:estados}
\end{figure}


% ════════════════════════════════════════════════════════════════════════════
\section{Índices Optimizados}
% ════════════════════════════════════════════════════════════════════════════

Se crearon índices para optimizar las consultas más frecuentes del sistema:

\begin{table}[H]
\centering
\renewcommand{\arraystretch}{1.3}
\begin{tabularx}{\textwidth}{l l l X}
    \toprule
    \textbf{Índice} & \textbf{Tabla} & \textbf{Tipo} & \textbf{Propósito} \\
    \midrule
    \texttt{idx\_reservas\_fecha\_inicio} & reservas & B-tree & Búsquedas por rango de fecha. \\
    \texttt{idx\_reservas\_estado} & reservas & B-tree & Filtrado por estado. \\
    \texttt{idx\_reservas\_fecha\_estado} & reservas & B-tree compuesto & Consultas combinadas fecha + estado. \\
    \texttt{idx\_reservas\_usuario} & reservas & B-tree & Historial de reservas por estudiante. \\
    \texttt{idx\_reservas\_espacio} & reservas & B-tree & Reservas por espacio. \\
    \texttt{idx\_reservas\_espacio\_rango} & reservas & GiST & Consultas de disponibilidad (espacio + rango temporal). \\
    \texttt{idx\_auditoria\_reserva} & auditoría & B-tree & Trazabilidad por reserva. \\
    \texttt{idx\_usuarios\_codigo} & usuarios & B-tree & Búsqueda por código estudiantil. \\
    \bottomrule
\end{tabularx}
\caption{Índices de la base de datos.}
\end{table}


% ════════════════════════════════════════════════════════════════════════════
\section{Vistas y Consultas Analíticas}
% ════════════════════════════════════════════════════════════════════════════

\subsection{Función: \texttt{fn\_disponibilidad\_espacio(UUID, DATE)}}

Genera las 15 franjas horarias de 1 hora (06:00--07:00, 07:00--08:00, \ldots, 20:00--21:00) para un espacio en una fecha dada, indicando si cada franja está disponible u ocupada:

\begin{lstlisting}[caption={Función de disponibilidad de franjas horarias.}]
SELECT * FROM fn_disponibilidad_espacio(
    'b0000000-0000-4bbb-bbbb-000000000004'::uuid,  -- TCH-001
    '2026-09-15'::date
)
WHERE disponible = TRUE;
\end{lstlisting}

\subsection{Vista: \texttt{v\_disponibilidad\_espacios}}

Wrapper sobre la función anterior para uso directo en herramientas BI. Muestra la disponibilidad del día actual para todos los espacios activos.

\subsection{Vista: \texttt{v\_resumen\_reservas}}

Vista consolidada con JOINs de contexto que expone: nombre del estudiante, facultad, espacio, deporte, horario en hora local de Bogotá, duración en horas, día de la semana y estado.

\subsection{Vista analítica: \texttt{v\_ocupacion\_franjas\_por\_deporte}}

Calcula la tasa de ocupación por franja horaria segmentada por deporte, con ranking de demanda. Útil para Data Engineering y Business Intelligence:

\begin{lstlisting}[caption={Consulta de ocupación por franjas y deporte.}]
SELECT deporte,
       franja_horaria,
       total_reservas,
       porcentaje_ocupacion,
       ranking_demanda
FROM v_ocupacion_franjas_por_deporte;
\end{lstlisting}


% ════════════════════════════════════════════════════════════════════════════
\section{Estructura de Archivos SQL}
% ════════════════════════════════════════════════════════════════════════════

El código fuente se organiza en 5 scripts modulares más un archivo consolidado:

\begin{table}[H]
\centering
\renewcommand{\arraystretch}{1.3}
\begin{tabularx}{\textwidth}{c l X}
    \toprule
    \textbf{Orden} & \textbf{Archivo} & \textbf{Contenido} \\
    \midrule
    1 & \texttt{01\_extensions\_and\_enums.sql} & Extensiones (\texttt{uuid-ossp}, \texttt{btree\_gist}) y 4 tipos ENUM. \\
    2 & \texttt{02\_schema\_ddl.sql} & 6 tablas, FK, CHECK constraints, EXCLUDE GiST, 8 índices. \\
    3 & \texttt{03\_triggers\_and\_validations.sql} & 3 funciones trigger: validación, auditoría, auto-timestamp. \\
    4 & \texttt{04\_seed\_data.sql} & 5 disciplinas, 12 espacios, 3 estudiantes, 3 reservas de prueba. \\
    5 & \texttt{05\_views\_and\_queries.sql} & Función de disponibilidad, 3 vistas, consultas analíticas. \\
    \midrule
    $\star$ & \texttt{DB-UnalSpace.sql} & \textbf{Archivo consolidado} --- contiene todo para ejecución directa. \\
    \bottomrule
\end{tabularx}
\caption{Estructura de archivos del proyecto.}
\end{table}


% ════════════════════════════════════════════════════════════════════════════
\section{Resumen de Integridad Referencial}
% ════════════════════════════════════════════════════════════════════════════

\begin{table}[H]
\centering
\renewcommand{\arraystretch}{1.3}
\begin{tabularx}{\textwidth}{l l l l X}
    \toprule
    \textbf{FK} & \textbf{Origen} & \textbf{Destino} & \textbf{On Delete} & \textbf{Justificación} \\
    \midrule
    \texttt{usuario\_id}    & reservas            & usuarios              & RESTRICT & No eliminar usuarios con reservas. \\
    \texttt{espacio\_id}    & reservas            & espacios\_deportivos  & RESTRICT & No eliminar espacios con reservas. \\
    \texttt{disciplina\_id} & reservas            & disciplinas           & RESTRICT & No eliminar disciplinas referenciadas. \\
    \texttt{espacio\_id}    & espacio\_deportes   & espacios\_deportivos  & CASCADE  & Si se borra un espacio, borrar sus asociaciones. \\
    \texttt{disciplina\_id} & espacio\_deportes   & disciplinas           & CASCADE  & Si se borra una disciplina, borrar asociaciones. \\
    \texttt{reserva\_id}    & auditoria\_reservas & reservas              & CASCADE  & Si se borra una reserva, borrar su log. \\
    \bottomrule
\end{tabularx}
\caption{Mapa de llaves foráneas e integridad referencial.}
\end{table}


\end{document}
